% 以section形式将附录加入到目录中
\addcontentsline{toc}{section}{附\hspace{2em}录}
\phantomsection    % 创建目录超链接目标
\phantom{ 空一行 }
\vspace{-1em}    % 修正center自带空行
\begin{center}
	{\zihao{3}\heiti 附\qquad 录}
\end{center}

% 代码格式的设置
\lstset{
  basicstyle=\ttfamily\small,  % 设置基本字体
  frame= shadowbox,   % 添加带阴影的边框
  numbers=left,     % 行号显示在左侧
  numberstyle=\tiny\color{gray},  % 行号颜色为灰色，tiny字体
  keywordstyle=\color{blue},    % 关键字 for ,while ,if等，改为蓝色
  commentstyle=\color{green},   % 注释样式：绿色
  stringstyle=\color{red},   % 字符串显示样式：红色
  breaklines=true,     % 长代码自动换行
  showstringspaces=true  % 字符串中的空格不显示
}

\newcounter{appendix_counter}
\newcommand{\theappsub}{\arabic{appendix_counter}}  % 新定义编号引用命令
\crefname{appsub}{附录}{附录}    % 定义正文进行\ref时所使用的格式


% 以subsection形式加入到目录当中
\refstepcounter{appendix_counter}   % 建立/ref节点 , counter会自动+1
\addcontentsline{toc}{subsection}{\theappsub \hspace{1em} Matlab代码示例 }  % 添加到目录中
\phantomsection
\arabic{appendix_counter}\hspace{1em} {\heiti\zihao{-4} Matlab代码示例}  % 标题
	\lstinputlisting[language=Matlab, label={matlab_example}]
		{codes/matlabcode.m}
\vspace{1em}

\refstepcounter{appendix_counter} 
\addcontentsline{toc}{subsection}{\theappsub\hspace{1em} Python代码示例 } 
\phantomsection
\arabic{appendix_counter}\hspace{1em} {\heiti\zihao{-4} Python代码示例}
	\lstinputlisting[language=Python, label={python_example}]
		{codes/pythoncode.py}
